\section{例题与分析}

\begin{question}\begin{example}[铜导线中的漂移速度]\label{c1:example:1}
铜导线直径 $d=2.0\,\mathrm{mm}$，电流 $I=7.6\,\mathrm A$。自由电子数密度 $n=8.5\times10^{28}\,\mathrm{m^{-3}}$，$e=1.60\times10^{-19}\,\mathrm C$。求漂移速率；电子平均漂移通过 $1.0\,\mathrm m$ 需多久？
\end{example}
\typeout{ECFIG:fig:c1:example:1:s}\fig{fig:c1:example:1:s}{题设：铜导线中的漂移速度}{\begin{tikzpicture}[scale=.88]
\draw[world](0,0)rectangle(4,1.1);\draw[guide](1.6,0)--(1.6,1.1);\draw[flow](.3,.55)--(1.3,.55)node[midway,above]{$I$};\draw[dim](0,-.5)--(4,-.5)node[midway,below]{$l$};\node[left]at(0,.55){$A$};\node[right]at(4,.55){$B$};\node[above]at(2,1.1){$S,\rho$};
\end{tikzpicture}}
\end{question}
\begin{solution}\textbf{按电流箭头确定电子漂移方向和截面积}\par\nobreak
取 图\ref{fig:c1:example:1:s} 中 $A\to B$ 为电流正向。电子带负电，因此电子漂移方向是 $B\to A$，求速率时用 $e$ 的大小。圆截面积为
\[
 S=\pi(d/2)^2=\pi\times10^{-6}\,\mathrm{m^2}.
\]


\par\medskip\Needspace{4\baselineskip}\noindent\textbf{以截面电荷输运求漂移速率}\par\nobreak
由时间 $\dd t$ 内穿过截面的电子数 $nSv_{\rm d}\dd t$ 得 $I=neSv_{\rm d}$，所以
\[
 v_{\rm d}=\frac{7.6}{8.5\times10^{28}\times1.60\times10^{-19}\times\pi\times10^{-6}}
 =1.78\times10^{-4}\,\mathrm{m/s}.
\]
\typeout{ECFIG:fig:c1:example:1:p}\fig{fig:c1:example:1:p}{解答过程：铜导线中的漂移速度}{\begin{tikzpicture}[scale=.88]
\draw[world](0,0)--(4.5,0);\draw[flow](.5,.5)--(4,.5)node[midway,above]{$I$};\draw[flow](4,-.5)--(.5,-.5)node[midway,below]{$v_{\rm d}$};\draw[dim](0,-1.2)--(4.5,-1.2)node[midway,below]{$L=1.0\,\mathrm m$};
\end{tikzpicture}}



\par\medskip\Needspace{4\baselineskip}\noindent\textbf{计算平均净位移所需时间并区别电路响应}\par\nobreak
图\ref{fig:c1:example:1:p} 把电子位移和正电流方向分开画出。平均漂移通过距离 $L$ 的时间为 $t=L/v_{\rm d}=5.62\times10^3\,\mathrm s$，约 $94\,\mathrm{min}$。这不表示电路接通后要等如此长时间才有响应：电路响应依靠电磁作用的传播，漂移时间描述的是载流子的平均净位移。\end{solution}
\begin{question}\begin{example}[两种离子共同导电]\label{c1:example:2}
电解液中正、负一价离子数密度均为 $n$。沿 $+x$ 的电场为 $E$，两种离子迁移率分别为 $\mu_+,\mu_-$，速率满足 $v_\pm=\mu_\pm E$。截面 $S$ 垂直 $x$ 轴。求电流、电导率和每单位体积的电功率。
\end{example}
\typeout{ECFIG:fig:c1:example:2:s}\fig{fig:c1:example:2:s}{题设：两种离子共同导电}{\begin{tikzpicture}[scale=.88]
\draw[axis](-.2,0)--(4.5,0)node[right]{$x$};\draw[axis](0,-.2)--(0,2.4)node[above]{$y$};\draw[flow](.4,1.4)--(1.8,1.4)node[midway,above]{$+e,\ v_+$};\draw[flow](3.8,.6)--(2.4,.6)node[midway,above]{$-e,\ v_-$};\node at(3,2){$\bm E\parallel+x$};
\end{tikzpicture}}
\end{question}
\begin{solution}\textbf{给两类离子的电荷和速度分别带符号}\par\nobreak
图\ref{fig:c1:example:2:s} 中正离子向右，负离子向左。分别代入带符号的 $q$ 与速度，有
\[
 j_+=ne\mu_+E,\qquad
 j_-=n(-e)(-\mu_-E)=ne\mu_-E.
\]


\par\medskip\Needspace{4\baselineskip}\noindent\textbf{相加求电流密度和电导率}\par\nobreak
两项都为正，故 $j=ne(\mu_++\mu_-)E$，$I=jS$，而电导率为 $\sigma=j/E=ne(\mu_++\mu_-)$。
\typeout{ECFIG:fig:c1:example:2:p}\fig{fig:c1:example:2:p}{解答过程：两种离子共同导电}{\begin{tikzpicture}[scale=.88]
\draw[world](0,0)rectangle(4,1.7);\draw[flow](.4,1.2)--(3.6,1.2)node[midway,above]{$j_+$};\draw[flow](.4,.5)--(3.6,.5)node[midway,below]{$j_-$};
\end{tikzpicture}}



\par\medskip\Needspace{4\baselineskip}\noindent\textbf{按两类载流子的总电流求局部功率}\par\nobreak
由 图\ref{fig:c1:example:2:p} 中两股同向的电流密度，体积功率是
\[
 p_V=(j_++j_-)E=ne(\mu_++\mu_-)E^2.
\]
迁移率单位 $\mathrm{m^2/(V\,s)}$，因而所得 $\sigma$ 的单位是 $\mathrm{S/m}$。若误把两种离子的速率直接相减，会遗漏电荷符号，甚至在 $\mu_+=\mu_-$ 时错误地得到零电流。\end{solution}
\begin{question}\begin{example}[同轴圆柱间的导电材料]\label{c1:example:3}
内外金属圆柱同轴，半径 $a,b$，长度 $h\gg b$，中间填充电阻率 $\rho$ 的均匀材料。内圆柱电势比外圆柱高 $U$，忽略端部效应。求 $I,j(r),E(r)$ 及电阻，并求半径 $r$ 处的体积发热功率。
\end{example}
\typeout{ECFIG:fig:c1:example:3:s}\fig{fig:c1:example:3:s}{题设：同轴圆柱间的导电材料}{\begin{tikzpicture}[scale=.88]
\draw[world](0,0)circle(.7);\draw[world](0,0)circle(1.6);\draw[dim](0,0)--(.7,0)node[midway,above]{$a$};\draw[dim](0,0)--(-1.13,-1.13)node[midway,left]{$b$};\node at(3,.5){$V(a)-V(b)=U$};\node at(3,-.3){长度 $h$};
\end{tikzpicture}}
\end{question}
\begin{solution}\textbf{由同轴对称性选径向串联薄层}\par\nobreak
在 图\ref{fig:c1:example:3:s} 中取半径 $r$、厚 $\dd r$ 的同轴薄层。由圆柱对称性，电流径向向外，各圆柱面通过同一总电流 $I$。图\ref{fig:c1:example:3:p} 标出有效面积是侧面 $2\pi rh$，不是圆盘面积。
\typeout{ECFIG:fig:c1:example:3:p}\fig{fig:c1:example:3:p}{解答过程：同轴圆柱间的导电材料}{\begin{tikzpicture}[scale=.88]
\draw[world](0,0)circle(.7);\draw[world](0,0)circle(1.6);\draw[guide](0,0)circle(1.1);\draw[guide](0,0)circle(1.25);\draw[flow](.8,0)--(1.5,0)node[right]{$I$};\node at(3.5,.5){$S(r)=2\pi rh$};\node at(3.5,-.3){$\dd R=\rho\,\dd r/(2\pi rh)$};
\end{tikzpicture}}



\par\medskip\Needspace{4\baselineskip}\noindent\textbf{由局域场积分求总电流和电阻}\par\nobreak
于是 $j=I/(2\pi rh)$，$E=\rho j$，薄层电压为 $\dd U=\rho I\,\dd r/(2\pi rh)$。积分得
\[
 U=\frac{\rho I}{2\pi h}\int_a^b\frac{\dd r}r
 =\frac{\rho I}{2\pi h}[\ln r]_a^b
=\frac{\rho I}{2\pi h}\ln\frac ba,\qquad
 I=\frac{2\pi hU}{\rho\ln(b/a)}.
\]


\par\medskip\Needspace{4\baselineskip}\noindent\textbf{求径向电流密度和体积发热分布}\par\nobreak
故 $R=\rho\ln(b/a)/(2\pi h)$，并有
\[
 j(r)=\frac{U}{\rho r\ln(b/a)},\quad
 E(r)=\frac{U}{r\ln(b/a)},\quad
 p_V(r)=\frac{U^2}{\rho r^2\ln^2(b/a)}.
\]


\par\medskip\Needspace{4\baselineskip}\noindent\textbf{积分总功率并检查窄间隙极限}\par\nobreak
发热在内侧更强。核验总功率时把体积元 $2\pi rh\dd r$ 同时代入：
\[
 P=\frac{2\pi hU^2}{\rho\ln^2(b/a)}\int_a^b\frac{\dd r}{r}
 =\frac{2\pi hU^2}{\rho\ln^2(b/a)}[\ln r]_a^b
 =\frac{2\pi hU^2}{\rho\ln(b/a)}=UI.
\]当 $b-a\ll a$ 时所得电阻回到平板间隙近似，模型自洽。\end{solution}
\begin{question}\begin{example}[电阻率渐变的细棒]\label{c1:example:4}
长 $l$、截面 $S$ 不变的细棒，电阻率为 $\rho(x)=\rho_0(1+\beta x/l)$，$0\le x\le l$，$\beta>-1$。两端电势为 $V(0)=U,V(l)=0$。求电阻、电场和电势分布。
\end{example}
\typeout{ECFIG:fig:c1:example:4:s}\fig{fig:c1:example:4:s}{题设：电阻率渐变的细棒}{\begin{tikzpicture}[scale=.88]
\draw[world](0,0)rectangle(4,1.1);\draw[guide](1.6,0)--(1.6,1.1);\draw[flow](.3,.55)--(1.3,.55)node[midway,above]{$I$};\draw[dim](0,-.5)--(4,-.5)node[midway,below]{$l$};\node[left]at(0,.55){$A$};\node[right]at(4,.55){$B$};\node[above]at(2,1.1){$S,\rho$};\node at(2,1.8){$\rho(x)=\rho_0(1+\beta x/l)$};
\end{tikzpicture}}
\end{question}
\begin{solution}\textbf{用稳恒条件确定同一电流}\par\nobreak
图\ref{fig:c1:example:4:s} 中没有分支，故 $I$ 与 $x$ 无关，电流密度也保持 $j=I/S$。电阻率却随 $x$ 变化，因而电场不能当作常数。

\par\medskip\Needspace{4\baselineskip}\noindent\textbf{积分渐变电阻率求电阻}\par\nobreak
微元串联给出
\[
 R=\int_0^l\frac{\rho_0}{S}\left(1+\frac{\beta x}l\right)\dd x
 =\frac{\rho_0}{S}\left[x+\frac{\beta x^2}{2l}\right]_0^l
=\frac{\rho_0l}{S}\left(1+\frac\beta2\right).
\]


\par\medskip\Needspace{4\baselineskip}\noindent\textbf{由局域欧姆关系求位置相关电场}\par\nobreak
由 $I=U/R$，电场为
\[
 E(x)=\rho(x)\frac IS
 =\frac{U(1+\beta x/l)}{l(1+\beta/2)}.
\]
\typeout{ECFIG:fig:c1:example:4:p}\fig{fig:c1:example:4:p}{取 $\beta=1$ 示意电场和电势趋势；推导保留一般 $\beta>-1$}{\begin{tikzpicture}

\begin{scope}[x=3.7cm,y=1.5cm]
\draw[axis](0,0)--(1.15,0)node[right]{$x/l$};\draw[axis](0,0)--(0,1.7)node[above]{$El/U$};
\draw[world,domain=0:1,samples=60]plot(\x,{2*(1+\x)/3});
\node[left]at(0,.666667){$2/3$};\draw[guide](1,0)--(1,1.333333);\node[below]at(1,0){$1$};
\end{scope}
\begin{scope}[xshift=6.5cm,x=3.7cm,y=1.5cm]
\draw[axis](0,0)--(1.15,0)node[right]{$x/l$};\draw[axis](0,0)--(0,1.7)node[above]{$V/U$};
\draw[world,domain=0:1,samples=80]plot(\x,{1-(\x+.5*(\x)^2)/1.5});
\node[left]at(0,1){$1$};\node[below]at(1,0){$1$};
\end{scope}

\end{tikzpicture}}



\par\medskip\Needspace{4\baselineskip}\noindent\textbf{积分求电势并核对两端边界}\par\nobreak
图\ref{fig:c1:example:4:p} 的微元电势降为 $E(x)\dd x$，所以从左端积分到 $x$，
\[
 V(x)=U-\int_0^xE(s)\dd s
 =U-\frac{U}{l(1+\beta/2)}
 \left[s+\frac{\beta s^2}{2l}\right]_0^x
=U-\frac{U}{1+\beta/2}\left(\frac xl+\frac{\beta x^2}{2l^2}\right).
\]
代入 $x=0,l$ 分别得到 $U,0$。当 $\beta=0$ 时，电场均匀、电势为直线；$\beta>0$ 时右侧材料电阻率大，同一电流对应更大的场，电势曲线下降得更快。若 $-1<\beta<0$，场强则沿 $+x$ 减小，电势下降逐渐变缓；两种情况下 $\rho(x)>0$ 都保证电势沿电流方向降低。\end{solution}
\begin{question}\begin{example}[半球电极与两脚间电势差]\label{c1:example:5}
均匀大地电阻率 $\rho$。半径 $a$ 的金属半球嵌入地面，向大地注入电流 $I$，回流电极非常远。地面两点距球心分别为 $r$ 和 $r+s$，$r>a$。求接地电阻、两点电势差，并判断何时可用 $Es$ 近似。
\end{example}
\typeout{ECFIG:fig:c1:example:5:s}\fig{fig:c1:example:5:s}{题设：半球电极与两脚间电势差}{\begin{tikzpicture}[scale=.88]
\draw[world](-2,0)--(4,0);\draw[world](-.5,0)arc(180:360:.5);\draw[guide](-1.4,0)arc(180:360:1.4);\fill(2,0)circle(1.5pt)node[above]{$A$};\fill(3,0)circle(1.5pt)node[above]{$B$};\draw[dim](0,.65)--(2,.65)node[midway,above]{$r$};\draw[dim](2,-.45)--(3,-.45)node[midway,below]{$s$};\node at(-1.5,-1.5){$\rho$};
\end{tikzpicture}}
\end{question}
\begin{solution}\textbf{利用地表边界选半球扩散面积}\par\nobreak
图\ref{fig:c1:example:5:s} 的地表是无电流穿出的边界，地下电流近似半球扩散。图\ref{fig:c1:example:5:p} 的半径 $x$ 处有效面积是 $2\pi x^2$，故
\[
 E(x)=\rho\frac{I}{2\pi x^2},\qquad
 V(x)=\frac{\rho I}{2\pi}\int_x^\infty u^{-2}\dd u
=\frac{\rho I}{2\pi}[-u^{-1}]_x^\infty=\frac{\rho I}{2\pi x}.
\]
\typeout{ECFIG:fig:c1:example:5:p}\fig{fig:c1:example:5:p}{解答过程：半球电极与两脚间电势差}{\begin{tikzpicture}[scale=.88]
\draw[world](-1.7,0)--(3.2,0);\foreach \r in {.55,1,1.5}{\draw[guide](-\r,0)arc(180:360:\r);}\foreach \a in {210,250,290,330}{\draw[flow](\a:.65)--(\a:1.4);}\node at(3,-.7){$j(x)2\pi x^2=I$};
\end{tikzpicture}}



\par\medskip\Needspace{4\baselineskip}\noindent\textbf{求电极电阻和两接触点电势差}\par\nobreak
电极的接地电阻 $R=V(a)/I=\rho/(2\pi a)$。两脚接触点间电势差为
\[
 U_{AB}=\frac{\rho I}{2\pi}\left(\frac1r-\frac1{r+s}\right)
 =\frac{\rho Is}{2\pi r(r+s)}.
\]


\par\medskip\Needspace{4\baselineskip}\noindent\textbf{精确比较局部匀强场近似的相对误差}\par\nobreak
若用 $E(r)s=\rho Is/(2\pi r^2)$，近似值与精确值之比为 $(r+s)/r=1+s/r$，相对高估 $s/r$。因此要求相对误差小于 $5\%$ 时，应有 $s/r<0.05$。该估计同时说明，远离电极时局部均匀场近似才可靠；不能把从球心到两脚的距离误当成两脚间距。\end{solution}
\begin{question}\begin{example}[注入与收集电极外侧的电场]\label{c1:example:6}
地表点电极 $P,Q$ 相距 $d$，分别向大地注入和收集电流 $I$。在 $Q$ 外侧同一直线上取点 $A$，$QA=r$，$r$、$d$ 远大于电极尺寸，大地电阻率为 $\rho$。求 $A$ 处电流密度、电场、电势，以及从 $A$ 向外走短距离 $s$ 的电势变化。
\end{example}
\typeout{ECFIG:fig:c1:example:6:s}\fig{fig:c1:example:6:s}{题设：注入与收集电极外侧的电场}{\begin{tikzpicture}[scale=.88]
\draw[world](-.5,0)--(6,0);\foreach \x/\t in {0/P,3/Q,4.6/A}{\fill(\x,0)circle(2pt)node[above]{$\t$};}\draw[flow](0,1)--(0,.2)node[midway,left]{$I$};\draw[flow](3,.2)--(3,1)node[midway,right]{$I$};\draw[dim](0,-.6)--(3,-.6)node[midway,below]{$d$};\draw[dim](3,-.6)--(4.6,-.6)node[midway,below]{$r$};
\end{tikzpicture}}
\end{question}
\begin{solution}\textbf{在电极外侧分别画注入和收集贡献}\par\nobreak
取 图\ref{fig:c1:example:6:s} 中 $P\to Q\to A$ 为 $+x$。注入电极在 $A$ 产生向右的电流密度，收集电极产生向左的电流密度。图\ref{fig:c1:example:6:p} 用两个相反箭头表示这两个贡献。
\typeout{ECFIG:fig:c1:example:6:p}\fig{fig:c1:example:6:p}{解答过程：注入与收集电极外侧的电场}{\begin{tikzpicture}[scale=.88]
\fill(0,0)circle(2pt)node[below]{$A$};\draw[flow](0,.4)--(1.6,.4)node[right]{$\bm j_P$};\draw[flow](0,-.4)--(-2,-.4)node[left]{$\bm j_Q$};\draw[axis](-2.4,-1.2)--(2,-1.2)node[right]{$x$};
\end{tikzpicture}}



\par\medskip\Needspace{4\baselineskip}\noindent\textbf{带符号叠加电流密度与电势}\par\nobreak
设无穷远电势为零，利用线性叠加，
\[
 j_x(A)=\frac{I}{2\pi(d+r)^2}-\frac{I}{2\pi r^2}<0,\qquad E_x=\rho j_x,
\]
\[
 V(A)=\frac{\rho I}{2\pi}\left(\frac1{d+r}-\frac1r\right)<0.
\]


\par\medskip\Needspace{4\baselineskip}\noindent\textbf{先求有限步长的精确电势变化}\par\nobreak
向外移动 $s$ 后的精确电势变化为
\[
 V(r+s)-V(r)=\frac{\rho I}{2\pi}
 \left[\frac1{d+r+s}-\frac1{r+s}-\frac1{d+r}+\frac1r\right].
\]


\par\medskip\Needspace{4\baselineskip}\noindent\textbf{检查小步长近似及电势变化符号}\par\nobreak
若 $s\ll r$ 且 $s\ll d+r$，可写成 $\Delta V\simeq-E_x(A)s>0$，即远离收集点时电势升高并趋近零。注意这里求的是“终点减起点”的电势变化；若求 $U_{A,A+s}$，符号应相反。点电极模型只能用于远处场，不能把 $r=0$ 代入计算总接触电阻。\end{solution}
\begin{question}\begin{example}[由宏观电阻率求电子的有效阻力]\label{c1:example:7}
金属中自由电子数密度为 $n$，电阻率为 $\rho$。电子以速率 $v$ 稳恒漂移。用宏观欧姆定律求单个电子所受平均阻力大小，并写出相应动量弛豫时间。
\end{example}
\typeout{ECFIG:fig:c1:example:7:s}\fig{fig:c1:example:7:s}{题设：由宏观电阻率求电子的有效阻力}{\begin{tikzpicture}[scale=.88]
\draw[world](0,0)rectangle(4,1.1);\draw[guide](1.6,0)--(1.6,1.1);\draw[flow](.3,.55)--(1.3,.55)node[midway,above]{$I$};\draw[dim](0,-.5)--(4,-.5)node[midway,below]{$l$};\node[left]at(0,.55){$A$};\node[right]at(4,.55){$B$};\node[above]at(2,1.1){$S,\rho$};\draw[flow](3,1.6)--(1,1.6)node[midway,above]{$v$};
\end{tikzpicture}}
\end{question}
\begin{solution}\textbf{从宏观电流密度反求平均驱动电场}\par\nobreak
图\ref{fig:c1:example:7:s} 中电流与电子漂移反向。先只取大小：$j=nev$，局域欧姆定律给 $E=\rho j=ne\rho v$。

\par\medskip\Needspace{4\baselineskip}\noindent\textbf{由平均受力平衡求阻力}\par\nobreak
漂移稳恒意味着平均加速度为零，故阻力大小等于电场力大小，
\[
 f_{\text{阻}}=eE=ne^2\rho v.
\]
\typeout{ECFIG:fig:c1:example:7:p}\fig{fig:c1:example:7:p}{解答过程：由宏观电阻率求电子的有效阻力}{\begin{tikzpicture}[scale=.88]
\fill(0,0)circle(2pt)node[above]{$-e$};\draw[flow](0,.4)--(-1.6,.4)node[left]{$\bm v_{\rm d}$};\draw[force](0,0)--(-1.6,0)node[left]{$-e\bm E$};\draw[force](0,-.4)--(1.6,-.4)node[right]{$\bm f_{\text{阻}}$};
\end{tikzpicture}}



\par\medskip\Needspace{4\baselineskip}\noindent\textbf{写出阻力矢量并确定弛豫时间}\par\nobreak
图\ref{fig:c1:example:7:p} 中电子向左，电场向右，所以电场力向左而平均阻力向右。写成矢量形式为 $\bm f_{\text{阻}}=-b\bm v_{\rm d}$，其中 $b=ne^2\rho$。若把 $b=m/\tau$，则
\[
 \tau=\frac{m}{ne^2\rho}.
\]


\par\medskip\Needspace{4\baselineskip}\noindent\textbf{核对时间尺度、量纲和漂移速度含义}\par\nobreak
这说明增大电阻率对应更快的平均动量弛豫。量纲检查：$b$ 的单位为 $\mathrm{kg/s}$，所以 $bv$ 为力。推导中的 $v$ 是平均漂移速率；若代入热运动速率，既不符合 $j=nev$ 的定义，也不能描述稳恒输运。\end{solution}
\begin{question}\begin{example}[金属棒突然停止后的简化漂移]\label{c1:example:8}
在晶格突然停止后的一个短暂阶段，电子相对晶格的初速度为 $v_0$，有效阻力为 $-ne^2\rho v$。设其他力和边界效应可忽略。求 $v(t)$、相对总位移，并解释为何没有有限的严格停止时刻。
\end{example}
\typeout{ECFIG:fig:c1:example:8:s}\fig{fig:c1:example:8:s}{题设：金属棒突然停止后的简化漂移}{\begin{tikzpicture}[scale=.88]
\draw[world](0,0)rectangle(4,1);\draw[flow](1.2,.5)--(2.8,.5)node[midway,above]{$v_0$};\node at(2,-.6){$t=0$：晶格停止};
\end{tikzpicture}}
\end{question}
\begin{solution}\textbf{选正向并按给定初值解衰减方程}\par\nobreak
按 图\ref{fig:c1:example:8:s} 取电子初速度方向为正，记 $\tau=m/(ne^2\rho)$。牛顿第二定律为 $\dot v=-v/\tau$。分离变量并代入 $v(0)=v_0$，
\[
 \begin{aligned}
 \int_{v_0}^{v(t)}\frac{\dd u}{u}
 &=-\frac1\tau\int_0^t\dd s,\\
 [\ln u]_{v_0}^{v(t)}&=-t/\tau,\qquad
 \ln\frac{v(t)}{v_0}=-t/\tau,\\
 v(t)&=v_0\e^{-t/\tau}.
\end{aligned}
\]


\par\medskip\Needspace{4\baselineskip}\noindent\textbf{积分累计相对位移}\par\nobreak
再对速度积分，
\[
 x(t)-x(0)=\int_0^t v_0\e^{-s/\tau}\dd s
 =v_0[-\tau\e^{-s/\tau}]_0^t
=v_0\tau(1-\e^{-t/\tau}).
\]
\typeout{ECFIG:fig:c1:example:8:p}\fig{fig:c1:example:8:p}{解答过程：金属棒突然停止后的简化漂移}{\begin{tikzpicture}[scale=.88]
\draw[axis](-.2,0)--(5,0)node[right]{$t/\tau$};\draw[axis](0,-.2)--(0,2.4)node[above]{$v/v_0$};\draw[world,domain=0:4.5,samples=60]plot(\x,{2*exp(-\x)});\node[left]at(0,2){$1$};\node at(3,1.5){面积 $\propto v_0\tau$};
\end{tikzpicture}}



\par\medskip\Needspace{4\baselineskip}\noindent\textbf{用曲线面积解释渐近停止与有限位移}\par\nobreak
图\ref{fig:c1:example:8:p} 的速度曲线在任何有限 $t$ 都大于零，却有有限面积，所以总位移为 $v_0\tau=mv_0/(ne^2\rho)$。也可对动力学方程从 $0$ 到 $\infty$ 积分：$-mv_0=-ne^2\rho\int_0^\infty v\dd t$，得到同一结果。两种计算分别核验时间变化和总位移。题设已经限制了模型；实际有限导体的电荷重新分布可能建立附加电场，不能把这一简化结果直接当作全过程。\end{solution}
\begin{question}\begin{example}[重力作用下竖直金属棒的电势差]\label{c1:example:9}
一竖直金属棒长 $l$，上端 $A$、下端 $B$。采用固定正离子背景、忽略晶格形变的简化模型，电子质量 $m$，电荷 $-e$。达到无宏观电流的平衡后，求内部电场和 $V_A-V_B$。
\end{example}
\typeout{ECFIG:fig:c1:example:9:s}\fig{fig:c1:example:9:s}{题设：重力作用下竖直金属棒的电势差}{\begin{tikzpicture}[scale=.88]
\draw[world](0,0)rectangle(.7,3);\node[above]at(.35,3){$A$};\node[below]at(.35,0){$B$};\draw[dim](1.4,0)--(1.4,3)node[midway,right]{$l$};\draw[force](2.7,2.5)--(2.7,1.3)node[right]{$\bm g$};
\end{tikzpicture}}
\end{question}
\begin{solution}\textbf{明确下向坐标中重力与电场力符号}\par\nobreak
图\ref{fig:c1:example:9:s} 中上、下端必须和重力方向一起判断。取竖直向下为 $+y$，电子重力为 $+mg$，电场力为 $-eE_y$。无相对漂移的条件为
\[
 mg-eE_y=0,\qquad E_y=\frac{mg}{e}>0.
\]


\par\medskip\Needspace{4\baselineskip}\noindent\textbf{核对电子受力图中的平衡方向}\par\nobreak
因此场向下，电子所受电场力向上，二者平衡。
\typeout{ECFIG:fig:c1:example:9:p}\fig{fig:c1:example:9:p}{解答过程：重力作用下竖直金属棒的电势差}{\begin{tikzpicture}[scale=.88]
\fill(0,1.5)circle(2pt)node[right]{$-e$};\draw[force](0,1.5)--(0,.3)node[right]{$m\bm g$};\draw[force](0,1.5)--(0,2.7)node[right]{$-e\bm E$};\draw[flow](2,2.7)--(2,.3)node[right]{$\bm E$};\node at(-.8,3){$A$};\node at(-.8,0){$B$};
\end{tikzpicture}}



\par\medskip\Needspace{4\baselineskip}\noindent\textbf{沿棒积分并检查电势高低及数量级}\par\nobreak
在 图\ref{fig:c1:example:9:p} 中沿 $A\to B$ 积分，$V_A-V_B=(mg/e)\int_0^l\dd y
=(mg/e)[y]_0^l=mgl/e$。上端电势较高，而不是根据电子重力方向直接猜测电势。若取 $l=1\,\mathrm m$，数量级约 $5.6\times10^{-11}\,\mathrm V$。此结论依赖固定晶格的假设，真实材料的压缩、弹性响应可能改变电荷分布，题中已把它们排除。\end{solution}
\begin{question}\begin{example}[旋转金属环的径向电势差]\label{c1:example:10}
金属圆环盘的内、外半径为 $a,b$，绕轴以角速度 $\omega$ 匀速转动。自由电子与晶格共同转动，电子质量 $m$、电荷 $-e$。忽略磁力及其他径向作用。求 $E_r(r)$ 与 $V(b)-V(a)$。
\end{example}
\typeout{ECFIG:fig:c1:example:10:s}\fig{fig:c1:example:10:s}{题设：旋转金属环的径向电势差}{\begin{tikzpicture}[scale=.88]
\draw[world](0,0)circle(.65);\draw[world](0,0)circle(1.65);\draw[dim](0,0)--(.65,0)node[midway,above]{$a$};\draw[dim](0,0)--(-1.16,-1.16)node[midway,left]{$b$};\draw[flow](1.17,1.17)arc(45:115:1.65)node[midway,above]{$\omega$};
\end{tikzpicture}}
\end{question}
\begin{solution}\textbf{列随晶格转动电子的径向运动方程}\par\nobreak
按 图\ref{fig:c1:example:10:s} 取径向向外为正。电子以角速度 $\omega$ 做圆周运动，径向加速度为 $-\omega^2r$。由电场力提供向心力，
\[
 -eE_r=-m\omega^2r,\qquad E_r(r)=\frac{m\omega^2r}{e}.
\]
\typeout{ECFIG:fig:c1:example:10:p}\fig{fig:c1:example:10:p}{解答过程：旋转金属环的径向电势差}{\begin{tikzpicture}[scale=.88]
\fill(0,0)circle(1.5pt)node[below]{$O$};\fill(1.6,0)circle(2pt)node[above]{$-e$};\draw[force](1.6,0)--(.5,0)node[midway,below]{$-eE_r$};\draw[flow](1.6,.6)--(2.8,.6)node[right]{$E_r$};\node at(1.6,-.9){$\bm a$ 指向 $O$};
\end{tikzpicture}}



\par\medskip\Needspace{4\baselineskip}\noindent\textbf{依据场向积分径向电势差}\par\nobreak
图\ref{fig:c1:example:10:p} 中电场向外，电子所受力向内，方向与向心加速度一致。由于 $E_r=-\dd V/\dd r$，有
\[
 V(b)-V(a)=-\int_a^bE_r(r)\dd r
 =-\frac{m\omega^2}{e}[r^2/2]_a^b
=-\frac{m\omega^2}{2e}(b^2-a^2).
\]


\par\medskip\Needspace{4\baselineskip}\noindent\textbf{检查窄环极限与忽略磁力的条件}\par\nobreak
因此外缘电势较低。当环很窄、宽度 $d=b-a\ll a$ 时，电势差近似为 $-(m\omega^2a/e)d$，等于用中间区域近似匀强场的结果。此题讨论惯性引起的静电分离，未引入外磁场，不能套用磁场中转动导体的动生电动势公式。\end{solution}
\begin{question}\begin{example}[串接粗细铜棒的电场与发热]\label{c1:example:11}
两根同材质铜棒长度均为 $l$，截面积分别为 $S$ 和 $2S$，串联接在恒定电压 $U$ 两端。忽略接合处端部效应，电阻率为 $\rho$。求各棒的电场、两端电压、单位体积发热功率之比。
\end{example}
\typeout{ECFIG:fig:c1:example:11:s}\fig{fig:c1:example:11:s}{题设：串接粗细铜棒的电场与发热}{\begin{tikzpicture}[scale=.88]
\draw[world](0,0)rectangle(2,.7);\draw[world](2,-.15)rectangle(4,.85);\node at(1,1.2){$S$};\node at(3,1.2){$2S$};\node at(0,-.5){$A$};\node at(2,-.6){$M$};\node at(4,-.5){$B$};\draw[flow](.2,1.8)--(1.5,1.8)node[midway,above]{$I$};
\end{tikzpicture}}
\end{question}
\begin{solution}\textbf{由串联同电流求两段电压}\par\nobreak
图\ref{fig:c1:example:11:s} 的 $M$ 为接合截面，稳恒电流要求两棒电流相同。两电阻为 $R_1=\rho l/S$、$R_2=\rho l/(2S)$，所以
\[
 I=\frac{U}{R_1+R_2}=\frac{2US}{3\rho l},\quad
 U_{AM}=IR_1=\frac{2U}{3},\quad U_{MB}=\frac U3.
\]
\typeout{ECFIG:fig:c1:example:11:p}\fig{fig:c1:example:11:p}{解答过程：串接粗细铜棒的电场与发热}{\begin{tikzpicture}[scale=.88]
\draw[world](0,0)rectangle(2,.7);\draw[world](2,-.15)rectangle(4,.85);\draw[flow](.1,.23)--(1.8,.23);\draw[flow](.1,.47)--(1.8,.47);\draw[flow](2.2,.35)--(3.9,.35);\node at(1,1.3){$E_1=2E_2$};\node at(3,-.7){$p_{V2}$};
\end{tikzpicture}}



\par\medskip\Needspace{4\baselineskip}\noindent\textbf{按各自面积求电流密度和电场}\par\nobreak
图\ref{fig:c1:example:11:p} 中细棒电流线较密，$j_1=I/S=2I/(2S)=2j_2$。由 $E=\rho j$，
\[
 E_1=\frac{2U}{3l},\qquad E_2=\frac{U}{3l}.
\]


\par\medskip\Needspace{4\baselineskip}\noindent\textbf{区分体积功率、总功率与初始温升速率}\par\nobreak
单位体积发热功率为 $\rho j^2$，所以 $p_{V1}:p_{V2}=4:1$；两棒总体积比为 $1:2$，总功率比则为 $2:1$。若隔热且单位体积热容量相同，初始温升速率比为 $4:1$。这三个比值对应不同物理量，不能互相替代。\end{solution}
\begin{question}\begin{example}[电源的功率分配与效率]\label{c1:example:12}
电源电动势 $\mathcal E$、内阻 $r$，接可变纯电阻 $R$。求电流、端电压、外电阻功率及供电效率；求外功率最大时的 $R$，并解释该点的效率。
\end{example}
\typeout{ECFIG:fig:c1:example:12:s}\fig{fig:c1:example:12:s}{题设：电源的功率分配与效率}{\begin{tikzpicture}[scale=.88]
\draw(0,0)to[battery1,invert,l={$\mathcal E$}](0,2.6)--(3.4,2.6)to[R,l={$R$}](3.4,0)to[R,l_={$r$}](0,0);\node[above right]at(3.4,2.6){$A$};\node[below right]at(3.4,0){$B$};\draw[flow](1,2.9)--(2.2,2.9)node[midway,above]{$I$};
\end{tikzpicture}}
\end{question}
\begin{solution}\textbf{先确定放电电流与端电压}\par\nobreak
图\ref{fig:c1:example:12:s} 中外电流从电源正端流出，电源内阻和外电阻串联，故
\[
 I=\frac{\mathcal E}{R+r},\qquad U=IR=\frac{\mathcal ER}{R+r}.
\]


\par\medskip\Needspace{4\baselineskip}\noindent\textbf{用能量分配求外功率和效率}\par\nobreak
电源每秒转换能量 $\mathcal EI$，外电阻获得 $P_R=I^2R=\mathcal E^2R/(R+r)^2$，效率为
\[
 \eta=\frac{P_R}{\mathcal EI}=\frac{R}{R+r}.
\]
\typeout{ECFIG:fig:c1:example:12:p}\fig{fig:c1:example:12:p}{归一化外功率随负载变化，峰值在 $R/r=1$，两侧极限均趋于零}{\begin{tikzpicture}

\begin{scope}[x=1cm,y=7cm]
\draw[axis](0,0)--(5.3,0)node[right]{$R/r$};\draw[axis](0,0)--(0,.36)node[above]{$P_R/(\mathcal E^2/r)$};
\draw[world,domain=0:5,samples=100]plot(\x,{\x/(1+\x)^2});
\draw[guide](1,0)--(1,.25)--(0,.25);\fill(1,.25)circle(1.3pt);
\node[below]at(1,0){$1$};\node[left]at(0,.25){$1/4$};
\end{scope}

\end{tikzpicture}}



\par\medskip\Needspace{4\baselineskip}\noindent\textbf{求导确定最大外功率}\par\nobreak
图\ref{fig:c1:example:12:p} 的功率曲线先增后减。求导得
\[
 \begin{aligned}
 \frac{\dd P_R}{\dd R}
 &=\mathcal E^2\left[(R+r)^{-2}-2R(R+r)^{-3}\right]\\
 &=\frac{\mathcal E^2[(R+r)-2R]}{(R+r)^3}
 =\frac{\mathcal E^2(r-R)}{(R+r)^3}.
\end{aligned}
\]
$R<r$ 时导数为正，$R>r$ 时为负，所以 $R=r$ 时 $P_{\max}=\mathcal E^2/(4r)$，此时 $\eta=1/2$。

\par\medskip\Needspace{4\baselineskip}\noindent\textbf{以极限区别高效率与大输出功率}\par\nobreak
当 $R\to\infty$，效率趋近 $1$，但电流和外功率均趋于零；最大功率和最大效率是不同的设计要求。\end{solution}
\begin{question}\begin{example}[角加速导线圈中的惯性电流]\label{c1:example:13}
细金属圆线圈半径 $R$，导线截面积 $S$，电阻率 $\rho$，自由电子数密度 $n$。线圈绕轴以角加速度 $\alpha$ 转动。电子质量 $m$、电荷 $-e$，平均阻力系数 $b=ne^2\rho$。忽略自感，且考查相对漂移已近似稳定的阶段。求相对晶格的电子速度与电流。
\end{example}
\typeout{ECFIG:fig:c1:example:13:s}\fig{fig:c1:example:13:s}{题设：角加速导线圈中的惯性电流}{\begin{tikzpicture}[scale=.88]
\draw[world](0,0)circle(1.4);\draw[flow](1.4,0)arc(0:70:1.4)node[midway,right]{$\alpha>0$};\draw[dim](0,0)--(-1.4,0)node[midway,above]{$R$};\node at(0,0){$O$};
\end{tikzpicture}}
\end{question}
\begin{solution}\textbf{取正切向并列电子所需切向加速度}\par\nobreak
取 图\ref{fig:c1:example:13:s} 的逆时针切向为正。相对漂移已稳定，电子的切向加速度近似等于晶格的 $R\alpha$。轴对称静电场沿整圈的环量为零，切向平均电场为零，因此电子所需的切向力由晶格的平均作用提供。

\par\medskip\Needspace{4\baselineskip}\noindent\textbf{以晶格作用力求相对漂移方向}\par\nobreak
以相对速度 $u$ 表示阻力，
\[
 mR\alpha=-bu,\qquad u=-\frac{mR\alpha}{ne^2\rho}.
\]
负号表示电子相对晶格顺时针漂移，常规电流为逆时针。
\typeout{ECFIG:fig:c1:example:13:p}\fig{fig:c1:example:13:p}{解答过程：角加速导线圈中的惯性电流}{\begin{tikzpicture}[scale=.88]
\draw[world](0,0)circle(1.4);\draw[flow](1.4,0)arc(0:65:1.4)node[midway,right]{$I$};\draw[flow](1.05,0)arc(0:-65:1.05)node[midway,right]{$u$};\node at(0,1.9){晶格角加速};
\end{tikzpicture}}



\par\medskip\Needspace{4\baselineskip}\noindent\textbf{由负电荷输运求常规电流}\par\nobreak
由 图\ref{fig:c1:example:13:p} 的方向关系，
\[
 I=neS|u|=\frac{mSR\alpha}{e\rho}.
\]


\par\medskip\Needspace{4\baselineskip}\noindent\textbf{用惯性驱动的环路积分独立复核}\par\nobreak
也可在线圈参考系中用电子惯性力 $-mR\alpha$ 等效为正电荷的驱动电场 $mR\alpha/e$，以弧长元 $R\dd\varphi$ 沿圈积分，得到
\[
 \mathcal E_{\rm eff}=\frac{mR\alpha}{e}\int_0^{2\pi}R\dd\varphi
 =\frac{mR^2\alpha}{e}[\varphi]_0^{2\pi}
 =\frac{2\pi mR^2\alpha}{e};
\]再除以线圈电阻 $2\pi R\rho/S$，得同一电流。此处能量来自使线圈加速的机械作用，且忽略自感的前提须由实际时间尺度保证。\end{solution}
\begin{question}\begin{example}[径向不均匀材料的轴向导电]\label{c1:example:14}
长直圆柱导体半径 $a$、长度 $l\gg a$，电导率只随径向变化：$\sigma(r)=\sigma_0(1-r^2/(2a^2))$。两端为等势金属端面，相差电压 $U$。忽略端部效应，求总电流、电阻和半径 $r$ 内通过的电流。
\end{example}
\typeout{ECFIG:fig:c1:example:14:s}\fig{fig:c1:example:14:s}{题设：径向不均匀材料的轴向导电}{\begin{tikzpicture}[scale=.88]
\draw[world](0,0)ellipse(.6 and 1);\draw[world](0,1)--(4,1)(0,-1)--(4,-1);\draw[world](4,0)ellipse(.6 and 1);\draw[flow](1,0)--(3,0)node[midway,above]{$I$};\node at(2,1.5){$\sigma=\sigma(r)$};\node at(-.9,0){$U$};\node at(5,0){$0$};
\end{tikzpicture}}
\end{question}
\begin{solution}\textbf{按公共等势端面确定轴向通道并联}\par\nobreak
图\ref{fig:c1:example:14:s} 的不同半径处可视作跨接同一对等势端面的细通道，它们是并联关系。各通道轴向电场相同，为 $E=U/l$，而电流密度 $j(r)=\sigma(r)U/l$ 随半径变化。
\typeout{ECFIG:fig:c1:example:14:p}\fig{fig:c1:example:14:p}{解答过程：径向不均匀材料的轴向导电}{\begin{tikzpicture}[scale=.88]
\draw[world](0,0)circle(1.4);\draw[guide](0,0)circle(.8);\draw[guide](0,0)circle(.95);\draw[dim](0,0)--(.8,0)node[midway,above]{$r$};\node at(3.5,.5){$\dd S=2\pi r\dd r$};\node at(3.5,-.4){各通道电压均为 $U$};
\end{tikzpicture}}



\par\medskip\Needspace{4\baselineskip}\noindent\textbf{对圆环电导积分求总电流和电阻}\par\nobreak
图\ref{fig:c1:example:14:p} 中薄圆环截面积 $\dd S=2\pi r\dd r$，故
\begin{align*}
 I&=\frac{\sigma_0U}{l}\int_0^a\left(1-\frac{r^2}{2a^2}\right)2\pi r\dd r\\
 &=\frac{2\pi\sigma_0U}{l}\left[\frac{r^2}{2}-\frac{r^4}{8a^2}\right]_0^a
 =\frac{3\pi\sigma_0a^2U}{4l}.
\end{align*}
因此$R=U/I=4l/(3\pi\sigma_0a^2)$。


\par\medskip\Needspace{4\baselineskip}\noindent\textbf{改变积分上限求内圈电流并核对边界}\par\nobreak
半径 $r$ 内的电流为
\[
 \begin{aligned}
 I(<r)&=\frac{2\pi\sigma_0U}{l}
 \int_0^r\left(s-\frac{s^3}{2a^2}\right)\dd s\\
 &=\frac{2\pi\sigma_0U}{l}
 \left[\frac{s^2}{2}-\frac{s^4}{8a^2}\right]_0^r
 =\frac{\pi\sigma_0U}{l}\left(r^2-\frac{r^4}{4a^2}\right).
\end{aligned}
\]
令 $r=a$ 回到总电流，令 $r=0$ 得零。这里的同轴薄通道是并联，而径向导电的同轴薄层是串联；几何形状相似并不意味着微元连接方式相同。\end{solution}
\begin{question}\begin{example}[双槽电解中的净质量变化]\label{c1:example:15}
三板从左到右为银板 $A$、铜板 $B$、银板 $C$，$A$ 接电源正端、$C$ 接负端，中间铜板 $B$ 不接外线并把两槽溶液隔开。左槽含 $\mathrm{Ag}^+$，右槽含 $\mathrm{Cu}^{2+}$，两侧电极过程均按给定价态计量。铜板一侧析出银，另一侧铜进入溶液，净质量以 $7.6\,\mathrm{mg/s}$ 增加。银、铜摩尔质量为 $108,64\,\mathrm{g/mol}$，$N_A=6.0\times10^{23}\,\mathrm{mol^{-1}}$、$e=1.6\times10^{-19}\,\mathrm C$，若用电压表测量两外板间电压，其电流可忽略。求电源电流。
\end{example}
\typeout{ECFIG:fig:c1:example:15:s}\fig{fig:c1:example:15:s}{两外板接电源，中间铜板隔开两槽溶液；忽略电压表电流}{\begin{tikzpicture}

\draw[world](0,0)rectangle(6,2.3);
\foreach \x/\lab in {1/A,5/C}{\draw[very thick](\x,.3)--(\x,1.8);\node[above]at(\x,1.8){$\lab$};}
\draw[very thick](3,0)--(3,2.3);\node[above]at(3,2.35){$B$};
\draw(1,1.8)--(1,3.3)--(2,3.3);\draw(4,3.3)--(5,3.3)--(5,1.8);
\node[draw,minimum width=2cm]at(3,3.3){电源};\node[above]at(1.5,3.3){$+$};\node[above]at(4.5,3.3){$-$};
\node at(2,.6){$\mathrm{Ag}^+$};\node at(4,.6){$\mathrm{Cu}^{2+}$};
\draw[flow](1.3,1.2)--(2.7,1.2);\draw[flow](3.3,1.2)--(4.7,1.2);
\node[below]at(3,-.3){$B$ 左侧析银，右侧铜溶入溶液};

\end{tikzpicture}}
\end{question}
\begin{solution}\textbf{从两槽串联和反应价态确定电量计数}\par\nobreak
图\ref{fig:c1:example:15:s} 中两槽串联，通过同一电量。银离子还原每摩尔需 $F=N_Ae=9.6\times10^4\,\mathrm{C/mol}$，铜溶出每摩尔释放 $2F$。电流 $I$ 对中板每秒析银质量为 $108I/F$ 克，溶铜质量为 $64I/(2F)$ 克。
\typeout{ECFIG:fig:c1:example:15:p}\fig{fig:c1:example:15:p}{中板两侧的反应与电子平衡；左侧增加银，右侧减少铜}{\begin{tikzpicture}

\fill[gray!20](2.35,0)rectangle(3.65,2.1);\draw[world](2.35,0)rectangle(3.65,2.1);
\node[above]at(3,2.5){中间铜板 $B$};
\draw[flow,thick](3.55,1.1)--(2.45,1.1);
\node[above]at(3,1.5){电子流};
\node[left,align=right]at(1.9,.2){$\mathrm{Ag}^++e^-\to\mathrm{Ag}$\\增加 $108I/F$ 克每秒};
\node[right,align=left]at(4.1,.2){$\mathrm{Cu}\to\mathrm{Cu}^{2+}+2e^-$\\减少 $64I/(2F)$ 克每秒};
\node[below]at(3,-.8){稳恒时中板不持续积累电荷，两槽电流相同};

\end{tikzpicture}}



\par\medskip\Needspace{4\baselineskip}\noindent\textbf{把析银与溶铜质量相减}\par\nobreak
按 图\ref{fig:c1:example:15:p} 净增加率不是单独析银率，而是两者相减：
\[
 \dot m=\frac IF(108-32)=\frac{76I}{96000}\,\mathrm{g/s}.
\]


\par\medskip\Needspace{4\baselineskip}\noindent\textbf{代入净增率并核对电极反应方向}\par\nobreak
代入 $\dot m=7.6\times10^{-3}\,\mathrm{g/s}$，得 $I=9.6\,\mathrm A$。若把铜按一价离子处理或忽略溶出质量，都会改变结果。槽电压并不直接进入这里的法拉第计量关系；本题由净质量增加率反求电流，电源接线只需保证图中的两侧反应方向。\end{solution}
